\documentclass[10pt]{article}
\usepackage[margin=0.65in]{geometry}
\usepackage[T1]{fontenc}
\usepackage[utf8]{inputenc}
\usepackage{lmodern}
\usepackage{enumitem}
\usepackage{amsmath,amssymb}
\usepackage{booktabs}
\usepackage{graphicx}
\usepackage{eso-pic}
\usepackage{xcolor}
\usepackage[hidelinks]{hyperref}
\setlength{\parindent}{0pt}
\setlength{\parskip}{0.45em}
\newcommand{\unitheader}[2]{%
  \vspace{0.35em}
  {\Large\bfseries Unit #1. #2\par}
  \vspace{0.2em}
  \hrule
  \vspace{0.55em}
}
\newcommand{\CodexWatermark}{%
  \AddToShipoutPictureBG{%
    \AtPageCenter{%
      \makebox(0,0){%
        \rotatebox{45}{%
          \textcolor[gray]{0.92}{\fontsize{22}{22}\selectfont\bfseries Winston Huang-Generated with Codex}%
        }%
      }%
    }%
  }%
}
\begin{document}
\CodexWatermark
{\centering
{\LARGE\bfseries AP Microeconomics Intensive Unit 1--6 Review\par}
\vspace{0.4em}
{\large Answer Key and Explanations\par}
\vspace{0.9em}}
\unitheader{1}{Basic Economic Concepts}
\begin{enumerate}[series=aps, leftmargin=*, itemsep=0.6em, topsep=0.2em, parsep=0pt]
\item \textbf{Answer: C}. Alder uses fewer labor hours per charger, while Briar uses fewer labor hours per bottle. Alder gives up only $\tfrac{1}{3}$ of a bottle per charger, compared with Briar's $\tfrac{3}{4}$, so Alder has the comparative advantage in chargers and Briar in bottles.
\item \textbf{Answer: B}. Alder specializes in chargers and produces $240/2=120$ chargers. Briar specializes in bottles and produces $240/4=60$ bottles.
\item \textbf{Answer: C}. The terms of trade must lie between the two opportunity costs of a charger: greater than Alder's $\tfrac{1}{3}$ bottle and less than Briar's $\tfrac{3}{4}$ bottle. Only $\tfrac{1}{2}$ bottle satisfies that condition.
\item \textbf{Answer: B}. Alder produces 120 chargers when specializing. Exporting 30 chargers leaves 90 chargers and brings in 15 bottles. Producing 90 chargers domestically would cost Alder 30 labor hours and leave only 10 bottles, so trade lets Alder consume beyond its own PPC.
\item \textbf{Answer: B}. Briar now has the absolute advantage in chargers because 1.5 hours is less than Alder's 2 hours. But Briar's opportunity cost of a charger is $1.5/4=0.375$ bottle, still above Alder's $2/6=\tfrac{1}{3}$, so Alder keeps the comparative advantage in chargers.
\end{enumerate}
\begin{enumerate}[resume*=aps, leftmargin=*, itemsep=0.6em, topsep=0.2em, parsep=0pt]
\item \textbf{Answer: D}. Moving from point C to point D raises drone output by 10 and lowers medical kits from 282 to 258, a sacrifice of 24 kits.
\item \textbf{Answer: B}. The losses in kits for each additional 10 drones are 6, 12, 24, 42, and 66. Because the sacrifice grows as production shifts toward drones, opportunity cost is increasing.
\item \textbf{Answer: B}. At 20 drones, the table shows a maximum of 282 kits. Producing only 240 kits means the economy is inside the frontier.
\item \textbf{Answer: B}. A technology improvement that affects only drone production increases the horizontal intercept while leaving the vertical intercept unchanged, so the PPC pivots outward around the medical-kit intercept.
\item \textbf{Answer: E}. Moving from point E to point F raises drones from 40 to 50 and lowers medical kits from 216 to 150, a decrease of 66 kits.
\end{enumerate}
\begin{enumerate}[resume*=aps, leftmargin=*, itemsep=0.6em, topsep=0.2em, parsep=0pt]
\item \textbf{Answer: B}. Comparing feasible bundles that use the \$24 budget, total utility is 93 for $(1,6)$, 108 for $(2,4)$, 105 for $(3,2)$, and 84 for $(4,0)$. The bundle with 2 guides and 4 coffees yields the highest total utility.
\item \textbf{Answer: C}. At that bundle, the next study guide gives 24 utils for \$6, or 4 utils per dollar. The sixth coffee gives 3 utils for \$3, or 1 util per dollar. Reallocating spending toward guides raises utility.
\item \textbf{Answer: D}. When income and all prices double by the same proportion, Lina's real purchasing power and relative prices are unchanged. Her optimal quantities stay the same.
\item \textbf{Answer: B}. A lower coffee price raises the marginal utility per dollar of coffee, making coffee relatively more attractive in utility maximization. Lina will generally substitute toward coffee.
\item \textbf{Answer: D}. Total utility is the sum of marginal utilities: $18+15+12+9=54$ utils.
\end{enumerate}
\begin{enumerate}[resume*=aps, leftmargin=*, itemsep=0.6em, topsep=0.2em, parsep=0pt]
\item \textbf{Answer: C}. Investment in capital goods and human capital expands productive capacity in the future, shifting the PPC outward more than current consumption does.
\item \textbf{Answer: C}. A positive statement is testable and descriptive. Whether a binding price ceiling creates a shortage is an empirical claim, not a value judgment.
\item \textbf{Answer: B}. A mixed market economy uses decentralized markets for most allocation while allowing government intervention for redistribution, externalities, and public goods.
\item \textbf{Answer: C}. The \$500{,}000 is a sunk cost and should not affect the current choice. Because marginal cost exceeds marginal benefit, the project should be rejected.
\item \textbf{Answer: C}. The city should continue as long as marginal benefit is at least marginal cost. That is true through the third drill, but not for the fourth.
\end{enumerate}
\unitheader{2}{Supply, Demand, and Market Intervention}
\begin{enumerate}[resume*=aps, leftmargin=*, itemsep=0.6em, topsep=0.2em, parsep=0pt]
\item \textbf{Answer: C}. Equilibrium occurs where quantity demanded equals quantity supplied. That happens at a price of \$6 and a quantity of 60 units.
\item \textbf{Answer: B}. With a binding ceiling, the quantity exchanged is limited by the short side of the market. At \$5, suppliers offer only 42 units even though consumers want 84.
\item \textbf{Answer: C}. At \$7, quantity supplied is 78 while quantity demanded is 60. The excess supply is 18 units.
\item \textbf{Answer: B}. A rise in income increases demand for a normal good. With supply fixed, the rightward demand shift raises both equilibrium price and equilibrium quantity.
\item \textbf{Answer: C}. Better technology shifts supply right, and higher income shifts demand left for an inferior good. Both shifts put downward pressure on price, so price falls for sure, but quantity depends on the relative sizes of the shifts.
\end{enumerate}
\begin{enumerate}[resume*=aps, leftmargin=*, itemsep=0.6em, topsep=0.2em, parsep=0pt]
\item \textbf{Answer: B}. A positive income elasticity means the service is normal, and a value greater than 1 indicates a luxury. A positive cross-price elasticity means bus rides and ride-sharing are substitutes.
\item \textbf{Answer: C}. When demand is elastic, a higher price causes quantity demanded to fall proportionally more than price rises, so total revenue falls.
\item \textbf{Answer: B}. A cross-price elasticity of $+0.8$ means the percentage change in ride-sharing demand equals $0.8$ times the percentage change in bus fares. A 5\% fall in bus fares therefore lowers ride-sharing demand by 4\%.
\item \textbf{Answer: B}. With an income elasticity of $+2.0$, demand changes by twice the percentage change in income in the same direction. A 3\% fall in income lowers quantity demanded by 6\%.
\item \textbf{Answer: A}. Tax incidence depends on relative elasticities, not on statutory assignment. The less elastic side of the market bears the larger share of the tax burden.
\end{enumerate}
\begin{enumerate}[resume*=aps, leftmargin=*, itemsep=0.6em, topsep=0.2em, parsep=0pt]
\item \textbf{Answer: B}. Set demand equal to supply: $30-0.5Q = 6+0.25Q$. Solving gives $Q=32$ and $P=14$.
\item \textbf{Answer: D}. The demand intercept is \$30, so consumer surplus is $\tfrac{1}{2}(30-14)(32)=256$.
\item \textbf{Answer: D}. The tax shifts supply up to $P = 9 + 0.25Q$. Setting this equal to demand gives $30 - 0.5Q = 9 + 0.25Q$, so $Q=28$ and the consumer price is \$16.
\item \textbf{Answer: B}. After the tax, 28 units are sold. Tax revenue equals \$3 times 28, which is \$84.
\item \textbf{Answer: B}. The tax reduces quantity from 32 to 28, so the lost trades equal 4 units. Deadweight loss is $\tfrac{1}{2}(3)(4)=6$.
\end{enumerate}
\begin{enumerate}[resume*=aps, leftmargin=*, itemsep=0.6em, topsep=0.2em, parsep=0pt]
\item \textbf{Answer: D}. At \$15, domestic quantity demanded is $120-2(15)=90$ and domestic quantity supplied is $20+2(15)=50$. Imports are therefore 40 units.
\item \textbf{Answer: B}. The autarky price is where $120-2P = 20+2P$, which gives \$25. The world price is lower, so consumers gain from cheaper coffee and producers lose from the lower domestic price.
\item \textbf{Answer: C}. A small-country tariff raises the domestic price to the world price plus the tariff, so the new domestic price is \$20.
\item \textbf{Answer: B}. At a domestic price of \$20, quantity demanded is 80 and quantity supplied is 60, so imports are 20 units. Tariff revenue is \$5 times 20, or \$100.
\item \textbf{Answer: C}. A quota of 20 imports forces the same import quantity as the tariff case. The domestic price must therefore rise to the level at which domestic demand exceeds domestic supply by 20 units, which is \$20.
\end{enumerate}
\unitheader{3}{Production, Cost, and Perfect Competition}
\begin{enumerate}[resume*=aps, leftmargin=*, itemsep=0.6em, topsep=0.2em, parsep=0pt]
\item \textbf{Answer: B}. Marginal product is the change in total product. The fourth worker raises output from 48 to 58, so marginal product is 10 units.
\item \textbf{Answer: A}. Marginal product falls from 18 for the first worker to 16 for the second worker, so diminishing marginal returns begin with the second worker.
\item \textbf{Answer: C}. Variable cost is the wage times the number of workers. With 5 workers at \$36 each, total variable cost is \$180.
\item \textbf{Answer: C}. Average fixed cost equals fixed cost divided by output. At 48 units, AFC is \$60/48 = \$1.25.
\item \textbf{Answer: B}. To raise output from 48 to 58, the firm must hire the fourth worker for \$36 and gains 10 units of output. Marginal cost is therefore \$36/10 = \$3.60 per unit.
\end{enumerate}
\begin{enumerate}[resume*=aps, leftmargin=*, itemsep=0.6em, topsep=0.2em, parsep=0pt]
\item \textbf{Answer: C}. The marginal costs are 12, 8, 8, 12, 18, and 26. The firm produces where price equals marginal cost on the upward-sloping portion of MC, which is 5 units when price is \$18.
\item \textbf{Answer: B}. At 5 units, total revenue is \$18 times 5, or \$90. Total cost is \$88, so economic profit is \$2.
\item \textbf{Answer: B}. The shutdown price is the minimum AVC. Variable costs are 12, 20, 28, 40, 58, and 84, so AVC values are 12, 10, 9.33, 10, 11.6, and 14. The minimum is about \$9.33.
\item \textbf{Answer: C}. A change in fixed cost does not change marginal cost, so the MR=MC output rule is unchanged in the short run.
\item \textbf{Answer: C}. At a price of \$10, the firm should still produce as long as price covers AVC. The closest MR=MC rule is 3 units: price exceeds the third unit's MC of 8 but is below the fourth unit's MC of 12.
\end{enumerate}
\begin{enumerate}[resume*=aps, leftmargin=*, itemsep=0.6em, topsep=0.2em, parsep=0pt]
\item \textbf{Answer: B}. In a constant-cost industry, entry shifts supply right until price returns to the original minimum-ATC level. Industry output rises, but long-run price returns to where firms earn zero economic profit.
\item \textbf{Answer: C}. Allocative efficiency occurs where the value to consumers, measured by price, equals the marginal cost of the last unit produced.
\item \textbf{Answer: C}. Productive efficiency means producing at the lowest possible average total cost, which is the minimum point of the ATC curve.
\item \textbf{Answer: B}. Losses cause firms to exit. Market supply shifts left, raising price until the remaining firms again earn zero economic profit.
\item \textbf{Answer: C}. A single competitive firm can sell as much as it wants at the market price, but nothing at a higher price, so its demand curve is horizontal.
\end{enumerate}
\begin{enumerate}[resume*=aps, leftmargin=*, itemsep=0.6em, topsep=0.2em, parsep=0pt]
\item \textbf{Answer: A}. Higher market demand raises price in the short run, giving firms positive profit. Entry then drives price back down until long-run profit returns to zero.
\item \textbf{Answer: C}. In an increasing-cost industry, expansion raises input prices and shifts firms' cost curves upward, so the new long-run equilibrium price is higher than before.
\item \textbf{Answer: A}. As long as price covers average variable cost, the firm can pay its variable inputs and contribute something toward fixed cost, making continued production better than shutting down.
\item \textbf{Answer: B}. When labor becomes more productive, each unit of output requires less variable input, so marginal cost falls and the MC curve shifts downward.
\item \textbf{Answer: C}. If all firms face higher per-unit costs, some firms exit until price rises enough for the remaining firms to earn normal profit again. In long-run competitive equilibrium, persistent losses cannot remain.
\end{enumerate}
\unitheader{4}{Imperfect Competition}
\begin{enumerate}[resume*=aps, leftmargin=*, itemsep=0.6em, topsep=0.2em, parsep=0pt]
\item \textbf{Answer: B}. The monopolist chooses output where $MR=MC$. Setting $50-2Q = 10$ gives $Q=20$.
\item \textbf{Answer: D}. At the profit-maximizing quantity of 20, the demand curve gives the price: $P = 50 - 20 = 30$.
\item \textbf{Answer: C}. Profit equals $(P-ATC)\times Q = (30-18)\times 20 = 240$.
\item \textbf{Answer: D}. Allocative efficiency occurs where price equals marginal cost. Setting $50-Q = 10$ gives $Q=40$.
\item \textbf{Answer: C}. Lower marginal cost shifts the MC curve down, so the profit-maximizing output rises and the monopoly price falls along the demand curve.
\end{enumerate}
\begin{enumerate}[resume*=aps, leftmargin=*, itemsep=0.6em, topsep=0.2em, parsep=0pt]
\item \textbf{Answer: D}. A perfectly price-discriminating monopolist expands output until price equals marginal cost, so it produces the allocatively efficient quantity of 40 units.
\item \textbf{Answer: B}. Perfect price discrimination makes the firm sell every unit for which willingness to pay exceeds marginal cost, so quantity rises from 20 to 40 and deadweight loss disappears.
\item \textbf{Answer: A}. A perfectly discriminating monopolist charges each buyer exactly their willingness to pay, capturing the entire surplus for itself.
\item \textbf{Answer: C}. A seller practicing third-degree price discrimination charges a lower price in the market segment with more elastic demand. Students are typically more price sensitive than business travelers.
\item \textbf{Answer: D}. Price discrimination relies on differences in willingness to pay or elasticity. If all consumers were identical, there would be no gain from charging different prices.
\end{enumerate}
\begin{enumerate}[resume*=aps, leftmargin=*, itemsep=0.6em, topsep=0.2em, parsep=0pt]
\item \textbf{Answer: C}. Both monopolistic competition and perfect competition feature many firms and low barriers to entry, even though firms in monopolistic competition sell differentiated products.
\item \textbf{Answer: A}. Because products are differentiated, a monopolistically competitive firm has some market power and therefore faces a downward-sloping demand curve, like a monopoly.
\item \textbf{Answer: B}. Positive profit attracts entry of close substitutes. That reduces demand for the original firm's product and shifts its demand curve left until economic profit is driven to zero.
\item \textbf{Answer: C}. In long-run monopolistic competition, the demand curve is tangent to ATC at an output smaller than the minimum-ATC quantity. The firm therefore has excess capacity and is not productively efficient.
\item \textbf{Answer: B}. Greater differentiation reduces the closeness of substitutes and makes consumers less price responsive, giving the firm more pricing power.
\end{enumerate}
\begin{enumerate}[resume*=aps, leftmargin=*, itemsep=0.6em, topsep=0.2em, parsep=0pt]
\item \textbf{Answer: C}. If FizzCo chooses High Price, ColaCo earns 16 by choosing Low Price instead of 12. If FizzCo chooses Low Price, ColaCo earns 8 by choosing Low Price instead of 4. So Low Price is dominant.
\item \textbf{Answer: D}. Because Low Price is the dominant strategy for both firms, the Nash equilibrium is for both to choose Low Price and earn 8 each.
\item \textbf{Answer: A}. Collusion seeks to maximize joint profit. Total profit is 24 at (High Price, High Price), which exceeds the 16 total at (Low Price, Low Price).
\item \textbf{Answer: A}. If one firm expects the other to keep a high price, it can increase its own payoff from 12 to 16 by cutting price. That incentive to cheat destabilizes collusion.
\item \textbf{Answer: A}. Without enforceable collusion, each firm has an incentive to follow its dominant strategy. The likely result is the noncooperative low-price equilibrium.
\end{enumerate}
\unitheader{5}{Factor Markets}
\begin{enumerate}[resume*=aps, leftmargin=*, itemsep=0.6em, topsep=0.2em, parsep=0pt]
\item \textbf{Answer: C}. The third worker adds 6 units of output, and each unit sells for \$12. The worker's marginal revenue product is therefore \$72.
\item \textbf{Answer: C}. The workers' MRPs are \$120, \$96, \$72, \$48, and \$24. The firm hires labor up to the point where MRP is at least the wage, so it hires 3 workers.
\item \textbf{Answer: B}. A higher output price raises the marginal revenue product of each worker, shifting the firm's labor demand curve to the right.
\item \textbf{Answer: C}. With a wage of \$45, the firm hires labor up to the point where MRP is at least \$45. The fourth worker's MRP is \$48, while the fifth worker's is \$24, so the firm hires 4 workers.
\item \textbf{Answer: B}. As more workers are hired with other inputs fixed, each additional worker tends to add less output. This causes marginal revenue product to fall and gives labor demand a downward slope.
\end{enumerate}
\begin{enumerate}[resume*=aps, leftmargin=*, itemsep=0.6em, topsep=0.2em, parsep=0pt]
\item \textbf{Answer: C}. Cost minimization requires $MP_L/w = MP_K/r$. Here $24/12 = 2$ and $36/18 = 2$, so the current combination is cost minimizing.
\item \textbf{Answer: B}. When the wage rises, $MP_L/w$ falls below $MP_K/r$. The bakery should substitute away from labor and toward ovens.
\item \textbf{Answer: B}. An increase in $MP_K$ raises $MP_K/r$ above $MP_L/w$, so the bakery can lower cost by using more ovens and less labor until the ratios are equal again.
\item \textbf{Answer: C}. Cost minimization is governed by $MP_L/w = MP_K/r$. Output price affects profit-maximizing scale, but not the least-cost input ratio for producing a given output.
\item \textbf{Answer: E}. The marginal revenue product of labor is price times marginal product: \$4 times 24 loaves, or \$96. Because this exceeds the \$12 wage, the bakery would want more labor at the margin.
\end{enumerate}
\begin{enumerate}[resume*=aps, leftmargin=*, itemsep=0.6em, topsep=0.2em, parsep=0pt]
\item \textbf{Answer: B}. Labor demand is a derived demand. When demand for the product increases, the demand for the labor used to produce that product also rises.
\item \textbf{Answer: B}. Better training raises marginal product, which increases marginal revenue product and shifts labor demand to the right.
\item \textbf{Answer: C}. A binding minimum wage raises the wage above equilibrium, causing quantity of labor supplied to exceed quantity of labor demanded.
\item \textbf{Answer: B}. When labor demand is relatively inelastic, the wage increase more than offsets the reduction in employment, so total earnings rise.
\item \textbf{Answer: B}. When capital becomes more expensive and labor can substitute for capital, firms shift toward labor, increasing labor demand.
\end{enumerate}
\begin{enumerate}[resume*=aps, leftmargin=*, itemsep=0.6em, topsep=0.2em, parsep=0pt]
\item \textbf{Answer: A}. The monopsonist hires where $MFC = MRP$: $10 + 4L = 70 - 2L$, so $L=10$. The wage is then read from the supply curve: $w = 10 + 2(10) = 30$.
\item \textbf{Answer: C}. Competition sets labor supply equal to labor demand: $10 + 2L = 70 - 2L$, so $L=15$ and $w=40$.
\item \textbf{Answer: C}. A monopsonist restricts employment to push wages down. Relative to competition, both wage and employment are lower.
\item \textbf{Answer: B}. A minimum wage set at the competitive wage can make the firm's marginal factor cost flat at \$40 up to the relevant quantity. The firm then hires where \$40 equals MRP, which is 15 workers.
\item \textbf{Answer: B}. In a monopsony, attracting an extra worker requires offering a higher wage, and that higher wage must generally be paid to all workers. Thus the marginal cost of one more worker exceeds the wage itself.
\end{enumerate}
\unitheader{6}{Market Failure and the Role of Government}
\begin{enumerate}[resume*=aps, leftmargin=*, itemsep=0.6em, topsep=0.2em, parsep=0pt]
\item \textbf{Answer: C}. The unregulated market sets demand equal to marginal private cost: $50-Q = 10+Q$. Solving gives $Q=20$.
\item \textbf{Answer: B}. Social efficiency requires demand to equal marginal social cost: $50-Q = 20+Q$. Solving gives $Q=15$.
\item \textbf{Answer: B}. The tax must make firms face the external cost. Since $MSC - MPC = 10$ at every quantity, the corrective tax is \$10 per unit.
\item \textbf{Answer: B}. The market overproduces from 15 to 20 units. The wedge between MSC and demand at 20 units is \$10, so deadweight loss is $\tfrac{1}{2}(5)(10)=25$.
\item \textbf{Answer: B}. When MSC exceeds MPC, firms base output on private rather than social cost. That causes the market to produce too much output.
\end{enumerate}
\begin{enumerate}[resume*=aps, leftmargin=*, itemsep=0.6em, topsep=0.2em, parsep=0pt]
\item \textbf{Answer: C}. For a public good, individual marginal benefits are added vertically. Thus $MSB = (30-Q)+(24-Q)+(18-Q) = 72-3Q$.
\item \textbf{Answer: C}. Set marginal social benefit equal to marginal cost: $72 - 3Q = 36$. Solving yields $Q=12$.
\item \textbf{Answer: B}. Because the barrier is nonexcludable, people can benefit without paying. That weakens private willingness to pay and leads to underprovision.
\item \textbf{Answer: D}. Pure public goods are both nonrival and nonexcludable, which is why competitive markets tend to underprovide them.
\item \textbf{Answer: C}. A common resource is generally nonexcludable but rival. An open-access fishery fits that description, unlike pure public goods such as national defense.
\end{enumerate}
\begin{enumerate}[resume*=aps, leftmargin=*, itemsep=0.6em, topsep=0.2em, parsep=0pt]
\item \textbf{Answer: B}. Natural monopolies have declining ATC over the relevant range, so marginal cost lies below ATC. Pricing at MC can therefore cause losses unless the firm is subsidized.
\item \textbf{Answer: B}. ATC pricing lets the firm cover all costs and earn normal profit, but because price remains above marginal cost, output stays below the allocatively efficient level.
\item \textbf{Answer: C}. Antitrust policy aims to preserve competition by discouraging mergers, collusion, and other practices that create or protect excessive market power.
\item \textbf{Answer: A}. Tradable permits direct pollution reduction toward firms that can reduce emissions most cheaply. That equalizes marginal abatement costs and lowers total cost.
\item \textbf{Answer: C}. In a monopsony, a moderate minimum wage can flatten marginal factor cost over a range of employment, leading the firm to hire more workers at a higher wage.
\end{enumerate}
\begin{enumerate}[resume*=aps, leftmargin=*, itemsep=0.6em, topsep=0.2em, parsep=0pt]
\item \textbf{Answer: C}. Country B gives larger cumulative income shares to the lower quintiles at every point, so its Lorenz curve is closer to the line of equality and its distribution is more equal.
\item \textbf{Answer: B}. A more equal Lorenz curve implies a smaller area between the curve and the line of equality, so Country B has the lower Gini coefficient.
\item \textbf{Answer: C}. Progressive taxation takes a larger share from higher incomes, and transfers raise lower incomes, both of which compress the after-tax income distribution.
\item \textbf{Answer: B}. Lower-income households tend to spend a larger share of their income, so a sales tax takes a larger percentage of income from them than from higher-income households.
\item \textbf{Answer: C}. Education and training increase skills and productivity, which raise earning potential and can reduce inequality over time through human-capital accumulation.
\end{enumerate}
\end{document}
